\documentclass[11pt,reqno]{amsart}

\usepackage[T1]{fontenc}
\usepackage[utf8]{inputenc}
\usepackage{lmodern}
\usepackage{amsmath,amssymb,amsthm,mathtools}
\usepackage{booktabs}
\usepackage{enumitem}
\usepackage[margin=1.05in]{geometry}
\usepackage{microtype}
\usepackage{xcolor}
\PassOptionsToPackage{hyphens}{url}
\usepackage[colorlinks=true,breaklinks=true,linkcolor=blue!60!black,citecolor=blue!60!black,urlcolor=blue!60!black]{hyperref}

\numberwithin{equation}{section}

\newtheorem{theorem}{Theorem}[section]
\newtheorem{proposition}[theorem]{Proposition}
\newtheorem{lemma}[theorem]{Lemma}
\newtheorem{corollary}[theorem]{Corollary}
\theoremstyle{definition}
\newtheorem{definition}[theorem]{Definition}
\newtheorem{example}[theorem]{Example}
\theoremstyle{remark}
\newtheorem{remark}[theorem]{Remark}

\newcommand{\C}{\mathbb{C}}
\newcommand{\Q}{\mathbb{Q}}
\newcommand{\Z}{\mathbb{Z}}
\newcommand{\GL}{\operatorname{GL}}
\newcommand{\diag}{\operatorname{diag}}

\title[Orthogonal Binary Fourier Masks]{Degree Classification, Block Rigidity, and Labeled Enumeration\\for Orthogonal Binary Fourier Masks}

\author[D. Kuttymbek]{Daniyar Kuttymbek}
\address{Astana, Kazakhstan}
\email{daniyar.kuttymbek@gmail.com}

\author[A. Kuttymbek]{Adilkhan Kuttymbek}
\address{Astana, Kazakhstan}
\email{adilkhan.kuttymbek@gmail.com}

\author[A. Kolpenova]{Amina Kolpenova}
\address{Astana, Kazakhstan}
\email{aminakolpenova@gmail.com}

\author[M. Kottumbekov]{Mukhammadali Kottumbekov}
\address{Astana, Kazakhstan}
\email{mukhammadalikottumbek@gmail.com}

\author[D. Kuttymbek]{Dayana Kuttymbek}
\address{Astana, Kazakhstan}
\email{dayanamiu74@gmail.com}

\date{\today}
\subjclass[2020]{Primary 05B20, 15B10; Secondary 11R18, 42A99, 15A60}
\keywords{Discrete Fourier matrix, binary masks, complex Hadamard submatrices, unit weighing matrices, vanishing sums of roots of unity, sparse inverse approximation}

\begin{document}

\begin{abstract}
Fix a positive integer $N$ and a primitive complex $N$-th root of unity $\zeta_N$. An \emph{orthogonal binary Fourier mask} of degree $d$ is an $N\times N$ binary matrix $R\in\{0,1\}^{N\times N}$ whose entrywise masked Fourier matrix $D_{rc}=R_{rc}\zeta_N^{rc}$ ($0\le r,c<N$) satisfies $DD^*=dI_N$ (equivalently $D^*D=dI_N$ for $d>0$). Such matrices sit strictly between two well-studied regimes: Cartesian Hadamard submatrices of Fourier matrices (where all selected rows share one common column support) and unit weighing matrices (where nonzero phases are unrestricted on the unit circle).

We establish five structural and enumerative results in this fixed-phase arbitrary-support model:
\begin{enumerate}[label=(\arabic*),leftmargin=1.8em]
\item \emph{Complete odd-semiprime degree classification:} For every pair of distinct odd primes $p<q$ and $N=pq$, an orthogonal binary Fourier mask of degree $d$ exists if and only if $d\in\{0,1,p,q,pq\}$, with no prime-ratio restriction.
\item \emph{Sharp sparse-inverse Fourier approximation:} For the normalized unitary Fourier matrix $F_N$, the operator-norm infimum $\Delta(N,s)=\inf\{\|M-F_N\|_2 : M\in\GL_N(\C),\ \text{every row of }M^{-1}\text{ has }\le s\text{ nonzeros}\}$ equals its lower bound $\sqrt{1-s/N}$ if and only if an orthogonal binary Fourier mask of degree $d=N-s$ exists, and is strictly larger otherwise. In particular, at odd semiprime $N=pq$, the floor is sharp precisely at $s\in\{pq-1,pq-p,pq-q\}$.
\item \emph{Semiprime block rigidity and exact labeled counts:} For every distinct prime pair $p<q$ (including $p=2$) and $d\in\{p,q\}$, every orthogonal mask of degree $d$ partitions into disjoint $d\times d$ blocks $A\times B$ in which both $A,B$ are transversals modulo $d$ and at least one is a coset of the order-$d$ subgroup; this yields an exact closed-form count of all labeled masks.
\item \emph{Mixed prime-power rigidity:} For $N=mq^b$ ($q$ prime, $b\ge 1$, $1\le m<q$), every orthogonal mask of degree $q$ partitions into disjoint $q\times q$ blocks characterized by uniform $q$-adic difference valuations $r+s=b-1$ and $m\mid (x-x')(y-y')$.
\item \emph{Arbitrary-$N$ weight-$3$ classification:} For every integer $N\ge 1$, an orthogonal binary Fourier mask of degree $3$ exists if and only if $3\mid N$, and every such mask splits into disjoint $3\times 3$ blocks---ruling out the $4\times 4$ one-hole blocks that occur for free-phase unit weighing matrices.
\end{enumerate}
\end{abstract}

\maketitle

\section{Introduction and main results}
\label{sec:intro}

Let $N\ge 1$ be an integer, let $\zeta=\zeta_N$ be a primitive complex $N$-th root of unity, and let $F_N\in\C^{N\times N}$ be the normalized unitary Fourier matrix with entries $(F_N)_{rc}=N^{-1/2}\zeta^{rc}$ ($0\le r,c<N$). For an arbitrary binary matrix $R\in\{0,1\}^{N\times N}$, define the unnormalized entrywise masked Fourier matrix
\begin{equation}
\label{eq:mask-def}
D_{rc} = R_{rc}\zeta^{rc},\qquad 0\le r,c<N.
\end{equation}
We say that $R$ (or $D$) is an \emph{orthogonal binary Fourier mask of degree $d$} if either $DD^*=dI_N$ or $D^*D=dI_N$.

Diagonal entries force $d\in\{0,1,\ldots,N\}$. When $d=0$, squared norms force $D=0$ and $R=0$. When $d=N$, every entry of $R$ must be $1$. Whenever $d>0$, either square identity $DD^*=dI_N$ or $D^*D=dI_N$ implies that $D/\sqrt{d}$ is unitary, hence implies the other identity automatically; consequently every row and every column of $R$ contains exactly $d$ ones.

This problem bridges two classical lines of research. On the rigid side, \emph{complex Hadamard submatrices} of $F_N$ select a single Cartesian rectangle $J\times K$ with $|J|=|K|=d$ (equivalently, spectral pairs in $\Z/N\Z$). At semiprime orders $N=pq$, divisor-only sizes $d\in\{1,p,q,pq\}$ follow from Bond--Glickfield--Herr~\cite[Theorem~1]{BGH19} and independently from Malikiosis--Kolountzakis~\cite[Theorem~1.2]{MK17} on Fuglede's conjecture; see also Herr--Wiegand~\cite{HW21}. In contrast, an orthogonal binary Fourier mask $R$ imposes neither Cartesian support nor block structure: row orthogonality depends on the pair-varying intersections $S_r\cap S_s$, so submatrix theorems do not apply without a global structural reduction. On the flexible side, \emph{unit weighing matrices} $\mathrm{UW}(N,d)$ allow arbitrary unit-modulus entries in place of $\zeta^{rc}$; Best--Kharaghani--Ramp~\cite{BKR13} showed that $\mathrm{UW}(N,3)$ exists for every $N\notin\{1,2,5\}$ (e.g., $\mathrm{UW}(35,3)$ exists via $3\times 3$ and $4\times 4$ blocks), whereas our Theorem~\ref{thm:weight3} shows that fixed Fourier phases restrict $d=3$ to $3\mid N$.

A direct analytic motivation comes from sparse inverse approximation of the discrete Fourier transform. For $N\ge 2$ and an integer row-sparsity budget $1\le s\le N-1$, define
\begin{equation}
\label{eq:delta-def}
\Delta(N,s) = \inf\bigl\{\|M - F_N\|_2 : M\in\GL_N(\C),\ \text{every row of } M^{-1}\text{ has at most }s\text{ nonzeros}\bigr\}.
\end{equation}
Every admissible matrix satisfies the elementary lower bound $\Delta(N,s)\ge \sqrt{1-s/N}$. Our first main result links sharpness of this bound directly to orthogonal binary Fourier masks.

\begin{theorem}[Sharpness of the inverse-row spectral floor]
\label{thm:inverse-row}
Let $N\ge 2$ and $1\le s\le N-1$, and put $t=s/N$ and $a=\sqrt{1-t}$. Then
\[
\Delta(N,s) \ge \sqrt{1-\frac{s}{N}},
\]
with equality if and only if there exists an orthogonal binary Fourier mask $R$ of degree $d=N-s$. Moreover:
\begin{enumerate}[label=(\alph*),leftmargin=2em]
\item If such a mask $R$ exists, set $C=F_N\circ R$ and $A=F_N-C$; then $A$ is invertible, the matrix $M = t(A^*)^{-1}$ (with $M^{-1}=t^{-1}A^*$) attains $\|M-F_N\|_2=a$, its inverse has exactly $s$ nonzeros in every row and column, and every minimizer attaining $\|M-F_N\|_2=a$ has this form.
\item If no such mask exists, then the infimum is strictly separated from the floor: $\Delta(N,s)>\sqrt{1-s/N}$.
\end{enumerate}
\end{theorem}

For any divisor $d\mid N$, a standard divisor-block construction (Section~\ref{sec:prelim}) produces an orthogonal binary Fourier mask of degree $d$, giving $\Delta(N,N-d)=\sqrt{d/N}$ for every proper divisor $d\mid N$. Whether $d\mid N$ is necessary for composite $N$ is a delicate global compatibility question. We settle it completely for all odd semiprimes.

\begin{theorem}[Odd-semiprime degree classification]
\label{thm:odd-semiprime}
Let $p<q$ be distinct odd primes, let $N=pq$, and let $\zeta$ be any primitive complex $N$-th root of unity. There exists a binary $N\times N$ matrix $R$ such that $D_{rc}=R_{rc}\zeta^{rc}$ satisfies $DD^*=dI_N$ (or $D^*D=dI_N$) if and only if
\begin{equation}
\label{eq:semiprime-deg}
d \in \{0, 1, p, q, pq\}.
\end{equation}
Equivalently, for $1\le s\le pq-1$,
\[
\Delta(pq,s) = \sqrt{1-\frac{s}{pq}} \iff s \in \{pq-1,\; pq-p,\; pq-q\},
\]
with strict inequality $\Delta(pq,s)>\sqrt{1-s/(pq)}$ at every other sparsity budget $s$.
\end{theorem}

Theorem~\ref{thm:odd-semiprime} asserts \emph{existence} at each allowed degree (clearly not every binary matrix of degree $d\in\{1,p,q\}$ is orthogonal). At the nontrivial divisor degrees $d\in\{p,q\}$, we go further and classify \emph{all} orthogonal support patterns and count them exactly, including even semiprimes $N=2q$.

\begin{theorem}[Semiprime block rigidity and exact labeled counts]
\label{thm:pq-rigidity}
Let $p<q$ be distinct primes (including $p=2$), $N=pq$, $d\in\{p,q\}$, and $k=N/d$.
\begin{enumerate}[label=(\alph*),leftmargin=2em]
\item \emph{(Block rigidity)} Every orthogonal binary Fourier mask $R$ of degree $d$ decomposes into $k$ disjoint $d\times d$ blocks $A_i\times B_i$ ($1\le i\le k$), where $A_1,\ldots,A_k$ and $B_1,\ldots,B_k$ partition $\Z/N\Z$.
\item \emph{(Local block characterization)} A $d\times d$ block $A\times B$ has orthogonal Fourier rows and columns if and only if both $A$ and $B$ are transversals modulo $d$ and at least one of $A,B$ is a coset of the order-$d$ subgroup $k\Z/N\Z$.
\item \emph{(Exact labeled count)} The total number of labeled $N\times N$ orthogonal binary Fourier masks of degree $d$ is
\begin{equation}
\label{eq:labeled-count}
\mathcal{M}(pq,d) = \sum_{\substack{0\le a,b\le k\\ a+b\ge k}} P_a P_b\, (b)_{k-a}\, a!,
\end{equation}
where $(b)_h = b!/(b-h)!$ and
\begin{equation}
\label{eq:Pa-def}
P_a = \binom{k}{a}\sum_{j=0}^{k-a} (-1)^j \binom{k-a}{j}\bigl((k-a-j)!\bigr)^{d-1}.
\end{equation}
\end{enumerate}
\end{theorem}

A similar prime-degree rigidity holds across mixed prime-power orders $N=mq^b$ whenever $m<q$.

\begin{theorem}[Fixed-degree rigidity at $N=mq^b$]
\label{thm:mqb}
Let $N=mq^b$, where $q$ is prime, $b\ge 1$, and $1\le m<q$. Every orthogonal binary Fourier mask of degree $d=q$ decomposes into $N/q$ disjoint $q\times q$ blocks $A\times B$. Moreover, a $q\times q$ block $A\times B$ is valid if and only if there exist nonnegative integers $r,s$ with $r+s=b-1$ such that:
\begin{enumerate}[label=(\roman*),leftmargin=2em]
\item $v_q(x-x')=r$ for all distinct $x,x'\in A$, and $v_q(y-y')=s$ for all distinct $y,y'\in B$; and
\item $m\mid (x-x')(y-y')$ for all $x,x'\in A$ and $y,y'\in B$.
\end{enumerate}
\end{theorem}

Finally, at weight $d=3$, we resolve existence and block structure for \emph{arbitrary} $N\ge 1$.

\begin{theorem}[Weight-$3$ Fourier masks for arbitrary $N$]
\label{thm:weight3}
Let $N\ge 1$ be any integer. An orthogonal binary Fourier mask of degree $d=3$ exists if and only if $3\mid N$. Moreover, every orthogonal binary Fourier mask of degree $3$ decomposes into $N/3$ disjoint $3\times 3$ blocks.
\end{theorem}

The paper is structured as follows. Section~\ref{sec:prelim} records the divisor-block construction, CRT coordinates, and classical prime-cycle cancellation. Section~\ref{sec:dense} proves the low-range and dense-range exclusions ($d<q$ and $d>p+q-1$) at odd semiprimes. Section~\ref{sec:window} introduces our global cycle-propagation lemma and cross-row incidence capacity bound to close the entire narrow window $q<d\le q+p-1$, completing Theorem~\ref{thm:odd-semiprime}. Section~\ref{sec:inverse} proves Theorem~\ref{thm:inverse-row} together with its prime and codimension-two corollaries. Section~\ref{sec:rigidity} proves Theorems~\ref{thm:pq-rigidity}, \ref{thm:mqb}, and~\ref{thm:weight3}. Section~\ref{sec:discussion} compares these results with prior literature and states open problems.

\section{Divisor constructions, CRT coordinates, and two-prime cancellation}
\label{sec:prelim}

\subsection{Explicit divisor-block masks}
For any positive divisor $d\mid N$, set $m=N/d$ and partition $\{0,\ldots,N-1\}$ into $m$ blocks of size $d$:
\begin{equation}
\label{eq:divisor-blocks}
I_t = \{td+u : 0\le u<d\},\qquad
J_t = \{t+\ell m : 0\le \ell<d\},\qquad 0\le t<m.
\end{equation}
Define $R_{rc}=1$ if $(r,c)\in \bigcup_{t=0}^{m-1} (I_t\times J_t)$ (equivalently $\lfloor r/d\rfloor \equiv c\pmod m$), and $R_{rc}=0$ otherwise. Every row and column has $d$ ones, and on block $t$:
\[
\zeta^{(td+u)(t+\ell m)} = \zeta^{(td+u)t}(\zeta^m)^{u\ell}
\]
because $td\ell m = t\ell N$. Since $\zeta^m$ is a primitive $d$-th root of unity, each $d\times d$ block is a diagonal unitary times an unnormalized $d\times d$ Fourier matrix, whose row and column Grams equal $dI_d$. Because the blocks occupy disjoint rows and columns, $DD^*=D^*D=dI_N$. Together with $R=0$ ($d=0$), this supplies sufficiency for every $d\mid N$ (in particular $d\in\{0,1,p,q,pq\}$ in Theorem~\ref{thm:odd-semiprime}).

\subsection{CRT coordinates and binary two-prime cancellation}
Let $p<q$ be distinct primes and $N=pq$. Choose the Chinese remainder idempotents $e_p = q(q^{-1}\bmod p)$ and $e_q = p(p^{-1}\bmod q)$ in $\Z/N\Z$. Writing original indices as $r = a e_p + b e_q$ and $c = u e_p + v e_q$ with $a,u\in\Z/p\Z$ and $b,v\in\Z/q\Z$, we have the exact identity
\begin{equation}
\label{eq:crt-phase}
\zeta^{rc} = \alpha^{au}\beta^{bv},\qquad \alpha = \zeta^{e_p},\quad \beta = \zeta^{e_q},
\end{equation}
where $\alpha$ is a primitive $p$-th root of unity and $\beta$ is a primitive $q$-th root of unity. No phases are modified: this is purely a coordinate relabeling of $\Z/N\Z\cong \Z/p\Z\times\Z/q\Z$.

For a binary mask $R$, denote by $S_r\subseteq \Z/p\Z\times\Z/q\Z$ the column support of row $r$ and by $T_c$ the row support of column $c$. Define the row groups, column $p$-fibers, and column $q$-lines by
\[
G_a = \{(a,b) : b\in\Z/q\Z\},\qquad
F_v = \{(u,v) : u\in\Z/p\Z\},\qquad
L_u = \{(u,v) : v\in\Z/q\Z\}.
\]
The orthogonality of two distinct rows $r=(a,b)$ and $s=(a',b')$ reads
\begin{equation}
\label{eq:row-orth}
\sum_{(u,v)\in S_r\cap S_s} \alpha^{(a-a')u}\beta^{(b-b')v} = 0.
\end{equation}
We recall the classical cyclotomic consequences (due to de Bruijn~\cite{deB53} and Lam--Leung~\cite[Theorem~3.3 and Corollary~3.4]{LL00}):
\begin{enumerate}[label=(C\arabic*),leftmargin=2.4em]
\item \emph{Same-group balance ($a=a'$, $b\neq b'$):} Equation~\eqref{eq:row-orth} becomes $\sum_{v\in\Z/q\Z} |S_r\cap S_s\cap F_v|\,\beta^{(b-b')v}=0$. Since $\Phi_q(X)=\sum_{v=0}^{q-1}X^v$ is irreducible over $\Q$, the fiber overlap count $|S_r\cap S_s\cap F_v|$ is independent of $v\in\Z/q\Z$. In particular, $|S_r\cap S_s|$ is a multiple of $q$.
\item \emph{Same-$b$ balance ($a\neq a'$, $b=b'$):} Symmetrically, the line overlap count $|S_r\cap S_s\cap L_u|$ is independent of $u\in\Z/p\Z$, and $|S_r\cap S_s|$ is a multiple of $p$.
\item \emph{Primitive binary two-prime cancellation ($a\neq a'$, $b\neq b'$):} By~\cite[Theorem~3.3]{LL00}, every vanishing sum of $pq$-th roots of unity with nonnegative integer weights is a nonnegative integer linear combination of full $p$-cycles and full $q$-cycles; equivalently, writing the indicator array $M_{uv}=\mathbf{1}_{(u,v)\in S_r\cap S_s}\in\{0,1\}$ in the cyclotomic ring $\Z[\alpha,\beta]$, vanishing of $\sum_{u,v}M_{uv}\alpha^{(a-a')u}\beta^{(b-b')v}$ forces $M_{uv} = A_u + B_v$ for rational functions of one coordinate. Since $M_{uv}\in\{0,1\}$, every pair of intersecting $p$- and $q$-cycles would create entry $2$, so $S_r\cap S_s$ must be either a union of complete $p$-fibers $F_v$ or a union of complete $q$-lines $L_u$ (including the empty set).
\end{enumerate}
Because $D^*D=dI_N$ holds simultaneously whenever $d>0$, the exact dual properties hold for column supports $T_c$.

\section{Low-range and dense-range exclusions at odd semiprimes}
\label{sec:dense}

Throughout Sections~\ref{sec:dense} and~\ref{sec:window}, let $p<q$ be distinct odd primes, $N=pq$, and let $R$ be an orthogonal binary Fourier mask of degree $d$.

\begin{proposition}[Low range $0<d<q$]
\label{prop:low-range}
If $0<d<q$, then $d\in\{1,p\}$.
\end{proposition}

\begin{proof}
By~(C1), distinct rows in a common group $G_a$ have overlap size divisible by $q$; since $d<q$, their supports are pairwise disjoint. Thus $qd = \sum_{r\in G_a}|S_r| \le pq$, giving $d\le p$. If $1<d<p$, then by~(C1)--(C3) every nonempty overlap of two distinct rows must have size at least $p>d$, so all $pq$ row supports are pairwise disjoint, giving $pqd\le pq$ and $d\le 1$, contradiction.
\end{proof}

\begin{proposition}[Dense range $p+q-1<d<pq$]
\label{prop:dense-range}
No orthogonal binary Fourier mask exists with $p+q-1<d<pq$.
\end{proposition}

\begin{proof}
Suppose $p+q-1<d<pq$. Fix any row $r=(a,b)$ and any column $c\in S_r$. At most $p+q-1$ rows agree with $r$ in at least one coordinate ($a'=a$ or $b'=b$). Since $|T_c|=d>p+q-1$, some row $s\in T_c$ differs from $r$ in both coordinates. By~(C3), $S_r\cap S_s$ is a union of complete $p$-fibers or complete $q$-lines containing $c$, so $S_r$ contains either the full $p$-fiber or the full $q$-line through $c$. Applying this to every $c\in S_r$ yields
\begin{equation}
\label{eq:dense-decomp}
S_r = (A_r\times\Z/q\Z)\cup(\Z/p\Z\times B_r),\qquad
d = \ell_r q + (p-\ell_r)f_r,
\end{equation}
where $A_r\subsetneq\Z/p\Z$ and $B_r\subsetneq\Z/q\Z$ are the maximal sets of complete $q$-lines and complete $p$-fibers contained in $S_r$, with $\ell_r=|A_r|<p$ and $f_r=|B_r|<q$.

For any two distinct rows $r,r'\in G_a$, the sets $B_r$ and $B_{r'}$ are disjoint: a common complete $p$-fiber $F_v$ gives $|S_r\cap S_{r'}\cap F_v|=p$, and~(C1) forces $|S_r\cap S_{r'}\cap F_w|=p$ for every $w\in\Z/q\Z$, implying $d\ge pq$, contradiction.

\smallskip
\noindent\emph{Case 1: $d=kq$ with $2\le k\le p-1$.} From~\eqref{eq:dense-decomp}, $(k-\ell_r)q=(p-\ell_r)f_r$. Since $1\le p-\ell_r\le p<q$, primality of $q$ forces $q\mid f_r$, hence $f_r=0$ and $\ell_r=k$. Every row support is therefore $A_r\times\Z/q\Z$ with $|A_r|=k$. Fixing $b\in\Z/q\Z$, for any two distinct rows $(a,b)$ and $(a',b)$, property~(C2) requires $|A_{(a,b)}\cap A_{(a',b)}|\in\{0,p\}$; since $k<p$, the $p$ sets $A_{(a,b)}$ ($a\in\Z/p\Z$) are pairwise disjoint $k$-subsets of $\Z/p\Z$, forcing $pk\le p$ and contradicting $k\ge 2$.

\smallskip
\noindent\emph{Case 2: $q\nmid d$.} Then $f_r\ge 1$ for every row $r$. Within each row group $G_a$, the $q$ nonempty sets $B_{(a,b)}$ ($b\in\Z/q\Z$) are pairwise disjoint subsets of $\Z/q\Z$, so every $B_{(a,b)}=\{\pi_a(b)\}$ is a singleton for a permutation $\pi_a$ of $\Z/q\Z$. Then $d = p + \ell_r(q-1)$, forcing $\ell_r=\ell$ to be constant across all rows, with $1\le\ell\le p-1$ since $p+q-1<d<pq$.

For distinct rows $r,r'\in G_a$, their two distinct special fibers each have overlap $\ell$, while any third fiber $F_w$ ($w\notin\{\pi_a(b),\pi_a(b')\}$) has overlap $|A_r\cap A_{r'}|$. By~(C1), $|A_r\cap A_{r'}|=\ell=|A_r|$, so $A_r=A_{r'}$. Thus every group $G_a$ shares a common $\ell$-subset $A_a\subset\Z/p\Z$. Let $h_u$ be the number of groups $a\in\Z/p\Z$ with $u\in A_a$. Counting the degree of any column $(u,v)$ gives
\[
d = h_u q + (p - h_u) = p + h_u(q-1),
\]
so $h_u=\ell$ for every $u\in\Z/p\Z$. Since $0<\ell<p$, the sets $A_a$ cannot all be equal (otherwise $h_u\in\{0,p\}$). Pick $a\neq a'$ with $A_a\neq A_{a'}$, and choose $r=(a,b)\in G_a$ and $s=(a',b')\in G_{a'}$ with $b'\neq b$ and $\pi_{a'}(b')\neq \pi_a(b)$ (at most two choices of $b'$ are forbidden, while $q\ge 5$). In the primitive overlap $S_r\cap S_s$, each of the two special fibers $\pi_a(b),\pi_{a'}(b')$ has overlap size $\ell\in\{1,\ldots,p-1\}$, whereas every other fiber has overlap $|A_a\cap A_{a'}|<\ell$. Thus $S_r\cap S_s$ is neither a union of full $p$-fibers nor a union of full $q$-lines, contradicting~(C3).
\end{proof}

\section{Cycle propagation and the narrow window $q<d\le q+p-1$}
\label{sec:window}

Propositions~\ref{prop:low-range} and~\ref{prop:dense-range} reduce the proof of Theorem~\ref{thm:odd-semiprime} to excluding the narrow window
\[
q < d \le q + p - 1.
\]
Note that every $d$ in this window satisfies $q<d<2q$, so by~(C1) any two distinct rows in the same group $G_a$ have fiber overlap $|S_r\cap S_s\cap F_v|\in\{0,1\}$ for all $v$.

\begin{lemma}[Cycle-propagation lemma]
\label{lem:cycle-prop}
If $q<d<pq$, then at least one row support $S_r$ contains a complete $q$-line $L_u$.
\end{lemma}

\begin{proof}
Suppose for contradiction that no row support contains a complete $q$-line. Call a supported entry $(r,c)$ with $r=(a,b)$ and $c=(u,v)$ \emph{deficient} if $F_v\not\subseteq S_r$. If some row $s\in T_c$ differed from $r$ in both coordinates, then~(C3) would force $S_r$ to contain either $L_u$ (impossible by assumption) or $F_v$ (impossible by deficiency). Therefore $T_c$ is contained in the coordinate cross
\[
X_r := \{(a,b') : b'\in\Z/q\Z\} \cup \{(a',b) : a'\in\Z/p\Z\}.
\]
Since $|T_c|=d>q$, column $c$ has at least one supporter $s=(a',b)$ with $a'\neq a$. For any other supporter $t=(a,b')\in G_a\cap T_c$ with $b'\neq b$, rows $t$ and $s$ differ in both coordinates and share $c$; since no $q$-line exists,~(C3) forces $F_v\subseteq S_t$ (and $F_v\subseteq S_s$).

There can be at most one such row $t\in G_a\setminus\{r\}$: if two distinct rows in $G_a$ both contained $F_v$, their overlap would have size $p$ in $F_v$, forcing overlap $pq>d$ by~(C1), impossible. Hence $|T_c\cap G_a|\le 2$ (consisting of $r$ and at most one $t$), while at most $p-1$ rows of $X_r$ lie outside $G_a$. Thus
\[
d = |T_c| \le 2 + (p-1) = p+1 < q+1 \le d,
\]
since $p<q$ are distinct odd primes (so $q\ge p+2$). This contradiction proves that no deficient entry exists.

Consequently, every row support $S_r$ is a union of complete $p$-fibers. Within any group $G_a$, distinct rows must use disjoint sets of $p$-fibers (again by~(C1) and $d<pq$). Since each row has $d>q>p$ entries, each row uses at least $2$ complete $p$-fibers, requiring at least $2q$ fibers in $\Z/q\Z$, contradiction.
\end{proof}

\begin{proposition}[Interior exclusion $q<d<q+p-1$]
\label{prop:interior}
No orthogonal binary Fourier mask exists with $q<d<q+p-1$.
\end{proposition}

\begin{proof}
Write $d=q+t$ with $1\le t\le p-2$. By Lemma~\ref{lem:cycle-prop}, some row $r=(a,b)$ contains a complete $q$-line $L_u$ plus $t$ extra columns outside $L_u$. Since $d<2q$, $L_u$ is the unique complete $q$-line in $S_r$. Because completing any $p$-fiber beyond its single baseline entry in $L_u$ requires $p-1>t$ extra columns, $S_r$ contains no complete $p$-fiber. Furthermore, $S_r$ meets at most $1+t\le p-1$ distinct $q$-lines, so $S_r$ is disjoint from at least one $q$-line $L_w$.

Pick any extra column $c\in S_r\setminus L_u$. Consider any row $s=(a',b')$ with $a'\neq a$:
\begin{itemize}[leftmargin=1.5em]
\item If $b'\neq b$, then $s$ cannot support $c$: by~(C3), a primitive overlap containing $c$ would force either the $q$-line or the $p$-fiber through $c$ into $S_r$, neither of which lies in $S_r$.
\item If $b'=b$, then~(C2) forces $|S_r\cap S_s\cap L_{u'}|$ to be constant over all $u'\in\Z/p\Z$. Since $S_r\cap L_w=\varnothing$, this constant is $0$, so $S_r\cap S_s=\varnothing$ and $s\notin T_c$.
\end{itemize}
Hence $T_c\subseteq G_a$, yielding $d=|T_c|\le q$, contradiction.
\end{proof}

\begin{proposition}[Boundary cross-row lemma]
\label{prop:cross}
If $d=q+p-1$, then every row $r$ containing a complete $q$-line $L_u$ must be an exact coordinate cross $S_r = L_u\cup F_v$ for some $v\in\Z/q\Z$.
\end{proposition}

\begin{proof}
Let $r=(a,b)$ contain its unique complete $q$-line $L_u$, and let $T=S_r\setminus L_u$ ($|T|=p-1$). Suppose for contradiction that $S_r$ contains no complete $p$-fiber. For every $c\in T$, no row differing from $r$ in both coordinates can support $c$ (by~(C3), since $S_r$ contains neither the $q$-line nor the $p$-fiber through $c$). Therefore $T_c\subseteq X_r$, where $|X_r|=q+p-1=d$. Regularity $|T_c|=d$ forces $T_c = X_r$ for \emph{every} $c\in T$.

Since $p\ge 3$, $|T|=p-1\ge 2$; pick two distinct columns $c_1=(u_1,v_1)$ and $c_2=(u_2,v_2)$ in $T$, and set $\delta_u=u_1-u_2$, $\delta_v=v_1-v_2$. Their column Gram entry on their common support $X_r$ is
\begin{equation}
\label{eq:cross-gram}
\sum_{(x,y)\in X_r} \alpha^{x\delta_u}\beta^{y\delta_v}
= \alpha^{a\delta_u}\beta^{b\delta_v}\bigl(q\,\mathbf{1}_{\delta_v=0} + p\,\mathbf{1}_{\delta_u=0} - 1\bigr),
\end{equation}
by inclusion--exclusion on the horizontal line $G_a$ and vertical fiber $\{(a',b)\}$ meeting at $(a,b)$. Because $c_1\neq c_2$, the parenthetical factor equals $-1$ (if $\delta_u,\delta_v\neq 0$), $p-1$ (if $\delta_u=0,\delta_v\neq 0$), or $q-1$ (if $\delta_u\neq 0,\delta_v=0$)---all nonzero---contradicting column orthogonality.

Thus $S_r$ contains a complete $p$-fiber $F_v$. Since $|F_v\setminus L_u|=p-1=|T|$, all $p-1$ extras lie in $F_v$, giving $S_r=L_u\cup F_v$.
\end{proof}

\begin{proof}[Proof of Theorem~\ref{thm:odd-semiprime}]
Sufficiency is given by~\eqref{eq:divisor-blocks} and $R=0$. By Propositions~\ref{prop:low-range}, \ref{prop:dense-range}, and~\ref{prop:interior}, it remains only to rule out $d=q+p-1$.

Assume $d=q+p-1$. By Lemma~\ref{lem:cycle-prop} and Proposition~\ref{prop:cross}, there exists a cross row $r=(a,b)$ with $S_r=L_u\cup F_v$. Put $T=F_v\setminus L_u$, so $|T|=p-1$. The total number of incidences between all $pq$ rows and the columns of $T$ must equal
\begin{equation}
\label{eq:inc-req}
\sum_{c\in T}|T_c| = (p-1)d = (p-1)(q+p-1).
\end{equation}
We upper-bound the same incidence sum row group by row group:
\begin{itemize}[leftmargin=1.5em]
\item \emph{The group $G_a$ of $r$:} Row $r$ contributes $p-1$ incidences to $T$. Since $d<2q$,~(C1) implies that every other row $t\in G_a\setminus\{r\}$ meets $S_r$ in at most $1$ point of $F_v$, hence in at most $1$ point of $T$. Thus $G_a$ contributes at most $(p-1)+(q-1)=q+p-2$ incidences.
\item \emph{Each other group $G_{a'}$ ($a'\neq a$):} First, if a row $s=(a',b')$ with $b'\neq b$ meets $T$, then~(C3) forces $F_v\subseteq S_s$ (as $S_r$ contains no $q$-line through any $c\in T$). At most one such row $s\in G_{a'}$ can exist, since two distinct rows in $G_{a'}$ containing $F_v$ would have overlap $pq>d$ by~(C1). Second, consider the unique row $s_0=(a',b)\in G_{a'}$ with the same second coordinate $b$. If $S_{s_0}\cap T\neq\varnothing$, then~(C2) forces $|S_r\cap S_{s_0}\cap L_{u'}|$ to be constant across all $u'\in\Z/p\Z$; since $|S_r\cap L_{u'}|=1$ for every $u'\neq u$, this constant must be $1$, forcing $T\subseteq S_{s_0}$.

Critically, a primitive supporter $s$ ($b'\neq b$) and the same-$b$ supporter $s_0$ cannot \emph{both} meet $T$: if both existed, they would be distinct rows of $G_{a'}$ sharing all $p-1\ge 2$ points of $T\subset F_v$, contradicting $|S_s\cap S_{s_0}\cap F_v|\le 1$ (since $d<2q$). Hence at most one row of $G_{a'}$ meets $T$, contributing at most $|T|=p-1$ incidences. Across all $p-1$ groups $G_{a'}$ ($a'\neq a$), the total contribution is at most $(p-1)^2$.
\end{itemize}
Summing both bounds and subtracting from~\eqref{eq:inc-req} yields the excess
\[
(p-1)(q+p-1) - \bigl(q+p-2 + (p-1)^2\bigr) = (p-2)(q-1) > 0,
\]
contradiction. This completes the degree classification~\eqref{eq:semiprime-deg}. The inverse-row consequence follows immediately from Theorem~\ref{thm:inverse-row} proved next.
\end{proof}

\begin{remark}[Why global column propagation is essential]
\label{rem:fano}
At $(p,q,d)=(11,23,33)$, there exists an explicit $23\times 253$ \emph{rectangular} binary Fourier-masked matrix $D_{\mathrm{rect}}$ (rows restricted to one group $G_a$, with two row-overlap clique components of sizes $7$ and $16$: five fibers carry Fano-plane lines on the $7$-row component, eighteen fibers share a singleton, and the $16$-row component uses disjoint baselines and $160$ private extras) satisfying the single-sided identity $D_{\mathrm{rect}}D_{\mathrm{rect}}^* = 33 I_{23}$. Thus a purely local row-group clique argument cannot eliminate $d=q+p-1$ when $q$ is small relative to $p^2$; the full square column-degree propagation in Lemma~\ref{lem:cycle-prop} and Propositions~\ref{prop:cross} is genuinely necessary.
\end{remark}

\section{Sharp sparse-inverse spectral approximation}
\label{sec:inverse}

We now prove Theorem~\ref{thm:inverse-row} and record its consequences for prime orders and codimension two.

\begin{proof}[Proof of Theorem~\ref{thm:inverse-row}]
Write $F=F_N$, $B=M^{-1}$, $E=M-F$, and let $b_i\in\C^{1\times N}$ denote row $i$ of $B$, with support $S_i$ of size $t_i\le s$ and norm $x_i=\|b_i\|_2>0$. Using $b_i E = e_i^T - b_i F$, $FF^*=I_N$, and Cauchy--Schwarz on $(b_i F)_i = \sum_{j\in S_i} (b_i)_j F_{ji}$ (where $|F_{ji}|=N^{-1/2}$):
\begin{align}
\|E\|_2^2 x_i^2 &\ge \|b_i E\|_2^2 = \|e_i^T - b_i F\|_2^2
= 1 + \|b_i\|_2^2 - 2\operatorname{Re}(b_i F)_i \nonumber\\
&\ge 1 + x_i^2 - 2\sqrt{t_i/N}\, x_i. \label{eq:one-row}
\end{align}
Dividing by $x_i^2$ and completing the square yields
\[
\|E\|_2^2 \ge \left(\frac{1}{x_i} - \sqrt{\frac{t_i}{N}}\right)^2 + 1 - \frac{t_i}{N} \ge 1 - \frac{s}{N} = a^2,
\]
proving $\Delta(N,s)\ge a$.

Suppose $\|E\|_2 = a$ is attained by $M$. Then equality holds throughout~\eqref{eq:one-row} for every $i$, forcing $t_i=s$, $x_i=t^{-1/2}$ (where $t=s/N$), and equality with positive real phase in Cauchy--Schwarz:
\[
(b_i)_j = t^{-1}\overline{F_{ji}}\,\mathbf{1}_{j\in S_i}.
\]
Setting $A_{ji} = F_{ji}\mathbf{1}_{j\in S_i}$, each column of $A$ has exactly $s$ nonzeros and $B = t^{-1}A^*$. Crucially, we now use all rows simultaneously: $H := a^2 I_N - EE^*$ is positive semidefinite, and equality in the first inequality of~\eqref{eq:one-row} states $b_i H b_i^* = 0$ for every $i$. Thus the positive semidefinite matrix $BHB^*$ has zero diagonal, hence $BHB^*=0$. Because $B$ is invertible, $EE^*=a^2 I_N$. Multiplying on left and right by $B$ and $B^*$ and substituting $BE=I_N-BF$, $FF^*=I_N$, and $a^2=1-t$ gives
\[
I_N - BF - F^*B^* + t BB^* = 0.
\]
Substituting $B=t^{-1}A^*$ and multiplying by $t$ yields $t I_N - A^*F - F^*A + A^*A = 0$. Setting the complementary mask $C = F - A = F\circ R$ (where $R\in\{0,1\}^{N\times N}$ has $N-s$ ones per column), we obtain
\begin{equation}
\label{eq:C-unitary}
C^*C = (F-A)^*(F-A) = I_N - F^*A - A^*F + A^*A = (1-t)I_N = a^2 I_N.
\end{equation}
Hence $D = \sqrt{N}\,C$ satisfies $D^*D = (N-s)I_N$, so $R$ is an orthogonal binary Fourier mask of degree $d=N-s$. Since $C/a$ is unitary, $CC^*=a^2 I_N$ as well, forcing every row of $R$ to have $N-s$ ones (and every row/column of $M^{-1}=t^{-1}A^*$ to have $s$ nonzeros).

Conversely, let $R$ be an orthogonal binary Fourier mask of degree $d=N-s$, and put $C=F\circ R$ and $A=F-C$. Since $\|C\|_2 = a < 1$ and $F$ is unitary, $\|Ax\|_2\ge (1-a)\|x\|_2$ for all $x\in\C^N$, so $A$ is invertible. Put $B=t^{-1}A^*$ and $M=B^{-1}=t(A^*)^{-1}$. Every row of $B$ has $s$ nonzeros, and reversing the algebra from~\eqref{eq:C-unitary} yields $BEE^*B^* = a^2 BB^*$, hence $EE^*=a^2 I_N$ and $\|M-F\|_2 = a$. (Notice also that the singular values of $M$ lie in $[1-a,1+a]$, so $\kappa_2(M)\le (1+a)/(1-a)$.)

Finally, to prove~(b), suppose $M_k\in\GL_N(\C)$ are admissible and $\|M_k-F\|_2\to a<1$. Then $\{M_k\}$ is bounded, so a subsequence converges to some $M_\infty\in\C^{N\times N}$. Since $\sigma_{\min}(M_k)\ge 1-\|M_k-F\|_2$ is bounded uniformly away from $0$ for large $k$, $M_\infty$ is invertible and $M_k^{-1}\to M_\infty^{-1}$. Because the set of matrices whose every row has at most $s$ nonzeros is closed (a finite union of coordinate subspaces), $M_\infty$ is admissible and attains $\|M_\infty-F\|_2=a$, yielding a mask $R$ by the necessity proof. Thus nonexistence of a mask implies $\Delta(N,s)>a$.
\end{proof}

\begin{corollary}[Divisors, codimension two, and prime dimensions]
\label{cor:easy-cases}
For $N\ge 2$:
\begin{enumerate}[label=(\arabic*),leftmargin=2em]
\item For every proper divisor $d\mid N$, $\Delta(N,N-d)=\sqrt{d/N}$. In particular, $\Delta(N,N-1)=1/\sqrt{N}$ for all $N\ge 2$.
\item For $N\ge 3$, $\Delta(N,N-2)=\sqrt{2/N}$ if $N$ is even, and $\Delta(N,N-2)>\sqrt{2/N}$ if $N$ is odd.
\item If $N=p$ is prime, then $\Delta(p,s)=\sqrt{1-s/p}$ if and only if $s=p-1$.
\end{enumerate}
\end{corollary}

\begin{proof}
Part~(1) follows from~\eqref{eq:divisor-blocks} and Theorem~\ref{thm:inverse-row}. For~(2) with $N$ odd, any orthogonal mask of degree $2$ has column overlap sums $\sum_{r\in T}\zeta^{r(k-j)}=0$ with $|T|\in\{0,1,2\}$; neither $|T|=1$ nor $|T|=2$ can vanish because $-1$ is not an $N$-th root of unity when $N$ is odd, so all $N$ column supports of size $2$ would be pairwise disjoint, impossible. For~(3), if $N=p$ is prime and $2\le d\le p-1$, then by irreducibility of $\Phi_p$ no nonempty proper subset of $p$-th roots sums to $0$, so all $p$ column supports of size $d$ must be pairwise disjoint, forcing $pd\le p$ and contradicting $d\ge 2$.
\end{proof}

\section{Block rigidity, labeled enumeration, and weight-$3$ masks}
\label{sec:rigidity}

\subsection{Semiprime rigidity and labeled counts (Theorem~\ref{thm:pq-rigidity})}

\begin{proof}[Proof of Theorem~\ref{thm:pq-rigidity}]
Let $p<q$ be distinct primes (now allowing $p=2$), $N=pq$, and $d\in\{p,q\}$.

\smallskip
\noindent\emph{Step 1: Block decomposition ($K_{d,d}$ components).}
At $d=p$, properties~(C1)--(C3) imply that every nonempty overlap of two row supports has size at least $p$, hence equals $p$; thus equal supports form isolated $K_{p,p}$ blocks (as column degree is $p$).

At $d=q$, proceed by induction on the number of remaining rows after removing all isolated $K_{q,q}$ blocks; suppose the remaining active square has size $tq$ with $t\ge 1$. Within any residue class $G_a$ of rows modulo $p$, property~(C1) implies that any nonempty overlap of two distinct rows has size $q$, so they share the exact same support. If no two active rows in any class $G_a$ shared a support, each of the $p$ classes would contain at most $t$ pairwise disjoint $q$-element supports (since the active square has only $tq$ columns), giving at most $tp$ active rows and contradicting $tp<tq$ (as $p<q$).

Hence some class $G_a$ contains two distinct active rows $r_1,r_2$ sharing a common support $B$ ($|B|=q$). Since $\sum_{c\in B}\zeta^{(r_1-r_2)c}=0$ and $p\mid (r_1-r_2)$ with $q\nmid (r_1-r_2)$, $B$ must be a transversal modulo $q$ (one element from each residue class mod $q$). Every row $r\in G_a$ meeting $B$ satisfies $S_r=B$ by~(C1). Now consider any row $s\notin G_a$ meeting $B$. Because $r_1\not\equiv r_2\pmod q$, at least one seed $r_i$ ($i\in\{1,2\}$) differs from $s$ modulo both $p$ and $q$. By~(C3), the nonempty primitive overlap $S_s\cap B$ is a union of complete $p$-fibers or complete $q$-lines; a complete $p$-fiber has $p$ elements with the same residue modulo $q$, impossible inside the mod-$q$ transversal $B$, so $S_s\cap B$ must be a full $q$-line and therefore exhausts $B$. Thus every row intersecting $B$ equals $B$. Column degree $q$ supplies exactly $q$ such rows, isolating a $K_{q,q}$ block $A\times B$ and completing the induction.

\smallskip
\noindent\emph{Step 2: Local block characterization.}
Let $A\times B$ be a $d\times d$ orthogonal block ($d\in\{p,q\}$). Suppose $A$ contained two distinct rows $x\equiv x'\pmod d$; then $x-x'$ is a nonzero multiple of $d$, so every term $\zeta^{(x-x')y}$ ($y\in B$) is an $(N/d)$-th root of unity. Because $N/d$ is prime, irreducibility of $\Phi_{N/d}$ implies that a vanishing sum of $d$ $(N/d)$-th roots of unity requires $(N/d)\mid d$, impossible for distinct primes $p,q$. Thus $A$ must be a transversal modulo $d$, and by $D^*D=dI_N$, $B$ must also be a transversal modulo $d$.

Now suppose $A$ is \emph{not} a coset of the order-$d$ subgroup $k\Z/N\Z$. Since $A$ is a transversal modulo $d$, two distinct elements $x,x'\in A$ must have $x\not\equiv x'\pmod k$, hence $\gcd(x-x',N)=1$. Their orthogonality $\sum_{y\in B}\zeta^{(x-x')y}=0$ is a primitive vanishing sum of $d$ roots of unity. By Lam--Leung~\cite[Theorem~3.3]{LL00}, any vanishing sum of $d$ primitive $pq$-th roots with $B$ a transversal modulo $d$ (which forbids any $k$-cycle: a $k$-cycle consists of $k$ elements congruent modulo $d$, impossible in a transversal mod $d$!) must be a single $d$-cycle, meaning all $y\in B$ share the same residue modulo $k$, i.e., $B$ is a coset of $k\Z/N\Z$. Conversely, if both $A,B$ are transversals modulo $d$ and at least one (say $B$) is a coset $y_0 + k\Z/N\Z$, then for any distinct $x,x'\in A$ we have $d\nmid (x-x')$ and
\[
\sum_{y\in B}\zeta^{(x-x')y} = \zeta^{(x-x')y_0}\sum_{j=0}^{d-1}(\zeta^k)^{(x-x')j} = 0,
\]
and symmetrically for columns (since every distinct $y,y'\in B$ has $y-y'\in k\Z/N\Z\setminus\{0\}$ and $A$ is a transversal mod $d$).

\smallskip
\noindent\emph{Step 3: Exact labeled count.}
A partition of $\Z/N\Z$ into $k=N/d$ transversals modulo $d$ is canonically indexed by the $k$ elements in residue class $0\bmod d$. Via the Chinese remainder bijection $\Z/N\Z \cong (\Z/k\Z)\times(\Z/d\Z)$, choosing the elements in the remaining $d-1$ residue classes modulo $d$ assigned to each block corresponds bijectively to choosing a $(d-1)$-tuple of permutations $(\pi_1,\ldots,\pi_{d-1})\in S_k^{d-1}$. A block $i\in\Z/k\Z$ is a coset of $k\Z/N\Z$ precisely when $i$ is a common fixed point of all $d-1$ permutations $\pi_1,\ldots,\pi_{d-1}$.

By inclusion--exclusion on the subset of common fixed points, the number of transversal partitions having \emph{exact} number $a$ of coset blocks is
\[
P_a = \binom{k}{a}\sum_{j=0}^{k-a}(-1)^j\binom{k-a}{j}\bigl((k-a-j)!\bigr)^{d-1},
\]
which is~\eqref{eq:Pa-def}. Given a row transversal partition with $a$ coset blocks and a column transversal partition with $b$ coset blocks, a pairing of the $k$ row blocks with the $k$ column blocks makes every matched block valid (by~(b)) if and only if each of the $k-a$ noncoset row blocks is matched to one of the $b$ coset column blocks. This is possible iff $b\ge k-a$ (i.e., $a+b\ge k$), in which case there are $(b)_{k-a}=b!/(b-k+a)!$ injections for the noncoset row blocks and $a!$ matchings for the remaining $a$ coset row blocks. Because distinct masks induce distinct connected $K_{d,d}$ components and hence distinct unordered block sets, summing over $a,b$ yields~\eqref{eq:labeled-count} with no overcounting.
\end{proof}

Table~\ref{tab:pq-masks} displays small exact values of~\eqref{eq:labeled-count}.

\begin{table}[ht]
\centering
\caption{Exact labeled counts $\mathcal{M}(pq,d)$ of orthogonal binary Fourier masks at semiprime orders.}
\label{tab:pq-masks}
\begin{tabular}{rrrr}
\toprule
Order $N=pq$ & Degree $d$ & $k=N/d$ & Labeled masks $\mathcal{M}(N,d)$ \\
\midrule
$6$ & $2$ & $3$ & $66$ \\
$6$ & $3$ & $2$ & $14$ \\
$10$ & $2$ & $5$ & $37{,}080$ \\
$10$ & $5$ & $2$ & $62$ \\
$15$ & $3$ & $5$ & $3{,}675{,}480$ \\
$15$ & $5$ & $3$ & $15{,}546$ \\
\bottomrule
\end{tabular}
\end{table}

\subsection{Mixed prime-power rigidity at $N=mq^b$ (Theorem~\ref{thm:mqb})}

\begin{proof}[Proof of Theorem~\ref{thm:mqb}]
Let $N=mq^b$ with $q$ prime, $b\ge 1$, and $1\le m<q$, and let $R$ be an orthogonal binary Fourier mask of degree $q$.

Remove any isolated $K_{q,q}$ blocks and suppose the remaining active square has size $tq$ ($t\ge 1$). Partition the active rows into $m$ classes according to their residue modulo $m$. Because $tq > tm$ (as $m<q$), some residue class modulo $m$ contains two distinct active rows $x,y$ whose supports intersect (if all supports in each class were pairwise disjoint, each class could own at most $t$ rows, totaling $\le tm < tq$).

Since $m\mid (x-y)$, the phases $\zeta^{(x-y)c}$ have $q$-power order $q^{b-r}$, where $r := v_q(x-y) < b$. A nonempty vanishing sum of at most $q$ $q$-power roots of unity must have length $q$ and form a single rotated $q$-cycle; hence $S_x = S_y =: B$ ($|B|=q$), and all distinct differences $c-c'$ ($c,c'\in B$) satisfy $v_q(c-c') = b-r-1$.

Now let $z$ be \emph{any} active row whose support meets $B$. By the ultrametric inequality on $v_q$, at least one difference $\delta\in\{z-x, z-y\}$ satisfies $v_q(\delta)\le r$ (since $v_q(x-y)=r$). Put $d' := b - v_q(\delta) \ge b - r \ge 1$ and write $\delta = q^{v_q(\delta)}\tilde{\delta}$ with $q\nmid \tilde{\delta}$. For all distinct $c,c'\in B$,
\[
v_q\bigl(\tilde{\delta}(c-c')\bigr) = b - r - 1 < b - r \le d',
\]
so the $q$-primary parts of $\zeta^{\delta c}$ ($c\in B$) take the form $\eta_0^{\tilde{\delta}c}$ for a primitive $q^{d'}$-th root of unity $\eta_0$, with pairwise distinct exponents $\tilde{\delta}c \pmod{q^{d'}}$.
Because $\gcd(m,q)=1$, the cyclotomic polynomial
\[
\Phi_{q^{d'}}(X) = \sum_{j=0}^{q-1} X^{j q^{d'-1}}
\]
remains irreducible over $\Q(\zeta_m)$ (by the cyclotomic degree formula $[\Q(\zeta_{mq^{d'}}):\Q(\zeta_m)]=\varphi(q^{d'})$). Therefore any nonzero polynomial $P(X)\in\Q(\zeta_m)[X]$ of degree $<q^{d'}$ vanishing at $\eta_0$ is a nonzero multiple of $\Phi_{q^{d'}}(X)$, so for every residue class modulo $q^{d'-1}$ appearing in $P$, all $q$ shifted powers $X^{a + j q^{d'-1}}$ ($0\le j<q$) must occur with the same nonzero coefficient in $\Q(\zeta_m)$. In particular, $P$ must have at least $q$ nonzero terms. Since $S_z\cap B$ is nonempty and yields such a polynomial with $|S_z\cap B|$ terms (all exponents distinct modulo $q^{d'}$), we must have $|S_z\cap B|=q$, i.e., $S_z=B$. By column degree $q$, exactly $q$ rows support $B$, isolating a $K_{q,q}$ block and completing the induction.

For the local block characterization of $A\times B$ ($|A|=|B|=q$): since $q>m$, by pigeonhole some distinct pair $x,x'\in A$ shares a residue modulo $m$. As above, orthogonality forces all distinct differences $y-y'$ ($y,y'\in B$) to have a single uniform $q$-adic valuation $s=b-1-v_q(x-x')$. Symmetrically, some pair in $B$ shares a residue modulo $m$, forcing all distinct differences $x-x'$ ($x,x'\in A$) to have a single uniform $q$-adic valuation $r$, with $r+s=b-1$. Now let $x,x'\in A$ be \emph{any} distinct pair, and set $\delta=x-x'$. Since $v_q(\delta)=r$, we have $d'=b-r=s+1$, and the $q$ values $\tilde{\delta}y\pmod{q^{s+1}}$ ($y\in B$) form a single complete congruence class modulo $q^s$, i.e., a full $\Phi_{q^{s+1}}$-progression. Because $\Phi_{q^{s+1}}(X)$ is irreducible over $\Q(\zeta_m)$, vanishing of $\sum_{y\in B}\omega_m^{\delta y}\eta_0^{\tilde{\delta}y}$ holds if and only if all $q$ coefficients $\omega_m^{\delta y}\in\Q(\zeta_m)$ are equal---which is equivalent to $m\mid \delta(y-y') = (x-x')(y-y')$ for all $y,y'\in B$. Conversely, conditions~(i) and~(ii) make every distinct row-pair inner product a constant phase times $\sum_{j=0}^{q-1}\zeta_q^j=0$, and symmetrically for columns.
\end{proof}

\subsection{Weight-$3$ Fourier masks for arbitrary $N$ (Theorem~\ref{thm:weight3})}

\begin{proof}[Proof of Theorem~\ref{thm:weight3}]
If $3\mid N$, the divisor-block construction~\eqref{eq:divisor-blocks} produces an orthogonal binary Fourier mask of degree $3$. Conversely, let $D_{rc}=R_{rc}\zeta_N^{rc}$ satisfy $DD^*=D^*D=3I_N$. Since $D$ is a unit weighing matrix of weight $3$, Best--Kharaghani--Ramp~\cite[Theorem~12]{BKR13} proves that $D$ (after row/column permutations) is a direct sum of $3\times 3$ full blocks and $4\times 4$ blocks having exactly one zero per row and column.

We prove that no $4\times 4$ one-hole block can occur with Fourier entries $\zeta_N^{rc}$. Suppose for contradiction that row labels $x_1,x_2,x_3,x_4\in\Z/N\Z$ and column labels $y_1,y_2,y_3,y_4\in\Z/N\Z$ form a $4\times 4$ block whose zeros lie on the diagonal $(x_i,y_i)$ ($1\le i\le 4$). For any partition of $\{1,2,3,4\}$ into two disjoint pairs $\{i,j\}$ and $\{k,\ell\}$, the overlap of rows $x_i$ and $x_j$ consists precisely of the two columns $\{y_k,y_\ell\}$. Orthogonality of rows $x_i,x_j$ reads
\[
\zeta_N^{(x_i-x_j)y_k} + \zeta_N^{(x_i-x_j)y_\ell} = 0
\iff
(x_i-x_j)(y_k-y_\ell) \equiv \frac{N}{2}\pmod N.
\]
If $N$ is odd, this is immediately impossible because $-1$ is not an $N$-th root of unity.

Assume now that $N$ is even, and set $e = v_2(N)\ge 1$. For every complementary pair of edges $\{i,j\}$ and $\{k,\ell\}$ of $K_4$, taking $2$-adic valuations yields
\begin{equation}
\label{eq:w3-val}
v_2(x_i-x_j) + v_2(y_k-y_\ell) = e - 1.
\end{equation}
Let $m_x := \min_{1\le i<j\le 4} v_2(x_i-x_j)$. Consider the graph $H_x$ on $\{1,2,3,4\}$ whose edges are the pairs achieving the minimum valuation $m_x$. Since $v_2$ is ultrametric (if $v_2(a-b)>m_x$ and $v_2(b-c)>m_x$, then $v_2(a-c)>m_x$), the complement of $H_x$ is a union of cliques, and because $2$ is prime, the residue classes of $x_i/2^{m_x}\pmod 2$ partition $\{1,2,3,4\}$ into exactly two nonempty parts of sizes either $2+2$ or $1+3$, with $H_x$ the complete bipartite graph between them.
\begin{itemize}[leftmargin=1.5em]
\item \emph{Case 1: $2+2$ split.} Then $H_x$ is a $4$-cycle; by~\eqref{eq:w3-val}, its complementary edges in $K_4$ form a $4$-cycle of edges $\{k,\ell\}$ on $\{y_1,y_2,y_3,y_4\}$ all having maximum valuation $M_y := e-1-m_x$. Writing $y_k \equiv y_\ell \pmod{2^{M_y}}$ along this $4$-cycle, transitivity forces the remaining two edges of $K_4$ (the two diagonals of the $4$-cycle) to satisfy $v_2(y_k-y_\ell)\ge M_y$. By maximality of $M_y$, all six column differences have valuation $M_y$; by~\eqref{eq:w3-val}, all six row differences have valuation $m_x$, impossible for four integers modulo $2^{m_x+1}$ (by pigeonhole, two of them agree mod $2^{m_x+1}$).
\item \emph{Case 2: $1+3$ split.} Say vertex $1$ is separated from $\{2,3,4\}$, so $v_2(x_1-x_j)=m_x$ for $j\in\{2,3,4\}$, while $v_2(x_j-x_k) > m_x$ for distinct $j,k\in\{2,3,4\}$. By~\eqref{eq:w3-val}, the complementary edges on $y$ satisfy $v_2(y_j-y_k) = M_y$ for all three pairs in $\{2,3,4\}$, producing a triangle of three integers $\{y_2,y_3,y_4\}$ whose pairwise differences all have the \emph{exact same} $2$-adic valuation $M_y$. Dividing by $2^{M_y}$ produces three integers with pairwise odd differences, impossible since there are only two parity classes modulo $2$.
\end{itemize}
Therefore no $4\times 4$ block exists. Every connected block of $D$ is $3\times 3$, forcing $3\mid N$.
\end{proof}

\section{Relation to prior work and open problems}
\label{sec:discussion}

Table~\ref{tab:compare} summarizes the structural boundary separating our fixed-phase binary Fourier-mask model from neighboring matrix families.

\begin{table}[ht]
\centering
\caption{Comparison of fixed-phase orthogonal binary Fourier masks with neighboring matrix classes.}
\label{tab:compare}
\begin{tabular}{llll}
\toprule
Model & Support structure & Nonzero phases & Semiprime / weight-$3$ behavior \\
\midrule
Hadamard submatrices of $F_N$~\cite{BGH19,HW21,MK17} & Single $d\times d$ rectangle $J\times K$ & Fixed $\zeta_N^{rc}$ & $d\in\{1,p,q,pq\}$ at $N=pq$ \\
Unit weighing $\mathrm{UW}(N,d)$~\cite{BKR13} & Arbitrary binary $R$ & Arbitrary on $\mathbb{T}$ & $\mathrm{UW}(N,3)$ exists $\forall N\notin\{1,2,5\}$ \\
Unitary digraphs~\cite{Sev03} & Arbitrary binary $R$ & Arbitrary $\C$ & Support criteria only \\
\textbf{Orthogonal Fourier masks (this paper)} & \textbf{Arbitrary binary $R$} & \textbf{Fixed $\zeta_N^{rc}$} & \textbf{$d\in\{0,1,p,q,pq\}$; $d=3\Leftrightarrow 3\mid N$} \\
\bottomrule
\end{tabular}
\end{table}

We also note the recent census of complex generalized weighing matrices~\cite{EOS26}, which treats bounded matrix orders with prescribed automorphism groups rather than arbitrary-support fixed-Fourier masks.

\subsection{Open problems}
\begin{enumerate}[label=(\arabic*),leftmargin=2em]
\item \emph{General composite orders:} Does every orthogonal binary Fourier mask of size $N$ satisfy $d\mid N$? The first open composite orders beyond odd semiprimes are $N=2pq$, $N=pqr$ (three distinct odd primes, starting at $N=105$), and $N=p^2 q$ (starting at $N=12, 18, 45$).
\item \emph{Even semiprimes $N=2q$:} Theorem~\ref{thm:pq-rigidity} resolves block rigidity and exact counts at $d\in\{2,q\}$, while Corollary~\ref{cor:easy-cases} handles $d=1,2$. Does an orthogonal binary Fourier mask of size $2q$ exist only for $d\in\{0,1,2,q,2q\}$?
\item \emph{Weight $4$ for arbitrary $N$:} Best--Kharaghani--Ramp~\cite{BKR13} classify $\mathrm{UW}(N,4)$; does fixed Fourier phase force $4\mid N$ (or allow $d=4$ when $2\mid N$)?
\item \emph{Quantitative spectral-norm gaps:} When no orthogonal binary Fourier mask of degree $N-s$ exists, Theorem~\ref{thm:inverse-row}(b) guarantees $\Delta(N,s)>\sqrt{1-s/N}$. Determine the exact value or asymptotic rate of the gap $\Delta(N,s)-\sqrt{1-s/N}$ (for instance, at prime orders $N=p$ and $1\le s\le p-2$).
\end{enumerate}

\section*{Data and code availability}
All verification artifacts are publicly available at
\begin{center}
\url{https://github.com/AgentSlop-Si/orthogonal-fourier-masks}
\end{center}
and consist of:
\begin{enumerate}[label=(\arabic*),leftmargin=2em]
\item A standalone Python verification script (\texttt{verify\_fourier\_masks.py}, standard library only) verifying the divisor-block masks, CRT identity, cross-character formula~\eqref{eq:cross-gram}, exact labeled mask counts at $N=6,10,15$ (including exhaustive brute-force verification at $N=6$ and $N=10,d=5$), the $N=mq^b$ local criterion at $N=18$, and the arbitrary-$N$ weight-$3$ obstruction up to $N=18$.
\item A zero-dependency Lean~4 formal verification file (\texttt{OrthogonalFourierMasks.lean}, verified in Lean~4.36.0 using core \texttt{Init} only, with no \texttt{sorry} and pure kernel reduction) machine-checking the divisor-block phase congruence of Section~\ref{sec:prelim}; the low-range and dense-range arithmetic bounds of Propositions~\ref{prop:low-range} and~\ref{prop:dense-range}; the cycle-propagation bound of Lemma~\ref{lem:cycle-prop}, cross-character nonvanishing of Proposition~\ref{prop:cross}, and global incidence excess $(p-2)(q-1)>0$ of Theorem~\ref{thm:odd-semiprime} for all odd primes $3\le p<q$; the odd-$N$ and even-$N$ parity/transitivity obstructions of Theorem~\ref{thm:weight3}; all six entries of Table~\ref{tab:pq-masks} from~\eqref{eq:labeled-count}--\eqref{eq:Pa-def}; an exact cyclotomic-ring ($\Z[\zeta_6]$) exhaustive arbitrary-mask enumeration at $N=6$ ($d\in\{2,3,4,5\}$); and the absence of $4\times 4$ one-hole Fourier blocks for all $4\le N\le 12$.
\end{enumerate}

\bibliographystyle{amsplain}
\bibliography{references}

\end{document}
